% !TEX root = ../main.tex
\section*{Chapter 5: Conclusion and Further Work}
\phantomsection \label{chap:conclusion} \setcounter{section}{5} \setcounter{subsection}{0}
\addcontentsline{toc}{section}{\textit{Chapter 5: Conclusion and Further Work}}

\subsection{Conclusion}

The broader question behind this project is what a traditional Linux system needs in order to host AI agents well, asked narrowly through MCP: which parts of a tool-invocation control plane need operating-system support, and which are fine in user space?
For the control-plane side the answer is a small set of invariants (semantic consistency between a registered tool and the one actually executed, identity binding between a request and its session, approval validity across the defer-approve-execute lifecycle, and state durability across gateway failure), each structurally hard to hold in a user-space gateway, because each has to survive the failure or compromise of the very process that enforces it.

linux-mcp is one concrete shape this split can take: a small kernel module holds registries, admission decisions, and the approval-ticket lifecycle behind a \texttt{CAP\_\allowbreak NET\_\allowbreak ADMIN}-gated Netlink interface, while \texttt{mcpd} handles manifests, semantic hashes, session binding, and forwarding.
\texttt{mcpd} is trusted to populate the kernel's view at initialisation and to forward Netlink frames faithfully thereafter, but its runtime logic (session state, ticket routing, request forwarding) is no longer load-bearing for individual request validation; a memory-safety compromise of the daemon itself remains out of scope, since \texttt{mcpd} still holds \texttt{CAP\_\allowbreak NET\_\allowbreak ADMIN}, and we discuss a minimal split of that capability in Section~\ref{sec:further-work}.
In the tested prototype, the kernel path enforced all specified control-plane invariants across the targeted attack set while both user-space baselines permitted meaningful bypass rates, with latency and throughput close to the user-space baseline throughout; a subsequent extended-attack rerun surfaced a missing peer-credential check in the stock module, closed by a six-line patch at no measurable latency cost.

The position is narrower than AIOS-style proposals that rebuild OS infrastructure around LLM abstractions, and less trusting than user-space frameworks that accept gateway correctness as given.
Whether AI agents will eventually warrant an operating system of their own remains an open question; in the meantime, a conventional Linux system extended with a few carefully scoped agent-aware mechanisms can already carry an MCP control plane that neither a user-space gateway nor existing Linux security primitives can match alone.

These control-plane benefits are not free at deployment time, and an honest account should name the costs that do not show up in the latency and throughput numbers.
Shipping a loadable kernel module introduces operational overhead that a user-space daemon avoids: a target-kernel version to track (the prototype is built against Linux 6.8 and compatibility across the active LTS range 5.15--6.12 has not been characterised), module-signing and loader-policy requirements on locked-down distributions, and a debugging path (\texttt{dmesg}, \texttt{kdb}, crash dumps) that most application engineers are unfamiliar with.
For deployments where user-space hardening is already sufficient, these overheads may outweigh the microsecond-scale gains reported here, and a realistic adoption path would likely require distribution-packaged modules, a documented kernel-version support matrix, and the \texttt{CAP\_\allowbreak NET\_\allowbreak ADMIN} split discussed in Section~\ref{sec:further-work} before the design is operationally comparable to pure user-space alternatives.

\subsection{Reflection}

The hardest part of this project was not writing the kernel code, but deciding what the kernel should not do.
For about the first six weeks, I treated AI-native OS support as a broad goal.
As a result, the design kept expanding: more in-kernel state, more policy hooks, and broader claims about what the kernel should provide for agents.
The project only became manageable when I narrowed it to four properties that a user-space gateway cannot reliably preserve on its own: tool identity, the approval-ticket lifecycle, session binding, and state durability across daemon failure.
Everything else had to stay out.
The rewrite after that decision was much faster, which in hindsight was a sign that the earlier version was not solving a harder problem, but a less clearly defined one.

I also spent roughly two weeks on the wrong kernel interface.
The first serious candidate was eBPF rather than Generic Netlink.
At first it looked attractive: pinned maps and \texttt{bpf\_\allowbreak timer} could have represented the required state.
The problem was the attachment point.
eBPF LSM hooks fire on kernel operations, not on parsed JSON-RPC messages, so an in-kernel arbitrator would have had to reconstruct MCP framing from raw socket traffic inside a verifier-bounded program on every invocation.
Plain LSM hooks had a similar problem: they run at the syscall boundary, where the notion of a tool does not yet exist.
Generic Netlink was the first interface that matched the granularity of the design, so I accepted the cost of a loadable kernel module.
That choice introduced its own problems.
Attribute-layout mismatches between the Python side and the kernel side produced malformed requests without compiler warnings or useful logs.
After losing nearly two days to one such mismatch, I started writing small per-command Python harnesses to construct and test each Netlink frame explicitly.

Another recurring temptation was to move more validation into the kernel.
Whenever I saw a check that a compromised \texttt{mcpd} might skip, the natural reaction was to duplicate it inside the module.
I made that mistake once with partial schema validation at registration time.
It worked until the schema changed; then a simple user-space update turned into a kernel rebuild and reload, while still duplicating logic that \texttt{mcpd} needed anyway.
Removing that code led to one of the more important design rules of the project: the kernel should verify \emph{membership}, but it should not compute \emph{content}.
This rule was not planned from the beginning.
It came from undoing a concrete mistake.
Without it, the module would have slowly grown into a second tool-execution engine inside the kernel, because every individual bug seemed to justify just one more kernel check.

The evaluation phase was the part that most changed how I think about systems experiments.
An early attack suite reported clean passes for the cross-UID hijack cases.
The result looked encouraging for several days, until I found that the kernel was not checking peer credentials on the \texttt{TOOL\_\allowbreak REQUEST} path at all.
The tests were passing because the harness was not exercising the path its name claimed to test.
The actual fix was small: compare the request peer uid with the uid bound to the session.
It added no measurable latency, but the uncomfortable lesson was that the earlier result had looked plausible enough to escape scrutiny.
A separate overload issue taught the same lesson in the opposite direction.
I initially blamed a kernel-side limit, but the real cause was \texttt{SOMAXCONN=128} saturation in the harness.
The habit I would carry into a future project is simple: a clean security-test pass is not enough.
I should first instrument the path and confirm that the intended check was actually reached.

\subsection{Further Work}
\label{sec:further-work}

The directions that follow are organised from short-range engineering items, which would consolidate the prototype into a deployment-ready system, to longer-range research questions for which the present design is closer to a starting point than a finished answer.

On the engineering side, several items remain inside the existing design envelope.
The kernel registries are protected by a single global lock that becomes a scalability bottleneck at higher concurrency and would benefit from RCU-protected lookup on the request path; hot reload is request-driven (\texttt{mcpd} recomputes a manifest-set signature on every request) rather than filesystem-driven, leaving a bounded staleness window between manifest edit and uptake that an \texttt{inotify}-driven path would close; cold-start reconciliation is unconditional rather than gated on an observable divergence signal; and the trust boundary around \texttt{mcpd} can be tightened by splitting \texttt{CAP\_\allowbreak NET\_\allowbreak ADMIN} into a minimal forwarder service that only relays pre-validated frames, which is consistent with the narrow-kernel principle and reduces the blast radius of a daemon compromise without widening the kernel interface.
The evaluation envelope has matching engineering items: bare-metal characterisation to eliminate hypervisor timing noise, $p_{99}$ under sustained concurrency past the clean-sweep ceiling (which requires repairing the \texttt{SOMAXCONN} saturation in the harness first), heterogeneous tool workloads with network-dependent execution times, and comparison against additional hardened user-space designs to separate benefits of kernel placement from those obtainable by user-space engineering alone.
None of these items changes the design narrative; they would close the gap between the prototype's measured behaviour and what a production deployment would need.

Beyond consolidation, four directions are closer to open research questions, where the present design is more naturally read as a starting point than a finished answer.

The first is a minimality argument for the kernel surface itself.
The present work picks four invariants by inspection (tool identity, approval-ticket lifecycle, session binding, state durability across daemon failure) and shows experimentally that a user-space gateway fails on each.
A stronger claim would establish a lower bound: under a stated threat model that admits memory-safety compromise of the gateway, prove that no purely user-space design can preserve all four properties simultaneously, and characterise the proposed kernel surface as not strictly larger than necessary.

The second is in-kernel, agent-attributed accounting and cooperative scheduling for tool execution.
The present prototype arbitrates whether a request runs but does not constrain how it runs once approved.
A natural next step is to extend the per-agent state already held in the kernel into a cgroup-style accounting domain that tracks tool-execution budgets (wall time, fan-out, tool-class quotas) and enforces them at request admission rather than after the fact, with multi-agent fairness guarantees backed by the same kernel state that already binds sessions.
This connects directly to the operating-system literature on resource-attribution containers, real-time scheduling, and BPF-CPU, and yields a falsifiable claim of the form ``two agents sharing a host see tool-execution latency that diverges by no more than $\Delta$ under contention''.

The third is attested arbitration.
At present, the kernel trusts \texttt{mcpd}'s view of the manifest set at registration time, and any remote party trusts the kernel's word on what was registered.
Anchoring the manifest-set hash in a TPM measurement and exposing a remote-attestation path that proves to an auditor which tool catalogue the kernel approved would close the gap between local enforcement and remote verifiability, which lines up with the broader confidential-computing direction in systems work.
The existing design already produces the necessary hashes; what is missing is a verifiable chain from kernel state through measured boot to a remote claim, together with a defensible answer for how that chain survives module unload, kernel live-patch, and approval-ticket reissue.

The fourth is multi-host federation.
The present design is single-machine, and any production deployment of LLM agents will span multiple hosts behind a fleet manager.
A federated variant in which kernel-anchored state is replicated across hosts via an attested gossip protocol, with a defined consistency model under partition (eventually-consistent registries, strongly-consistent approval tickets), would extend the narrow-kernel principle into a distributed-systems setting and address what is, today, the most obvious deployment-scale objection to the design.
Whether the kernel state needs to participate in the replication path or only sign for it, and whether approval-ticket linearisability requires consensus or can be reduced to leader-pinned issuance, are themselves the substantive research questions.

These four directions share a common shape: they take the current design's narrow kernel surface as a fixed boundary and ask what additional guarantees the user-space side can be given, rather than what additional logic the kernel can absorb.
That orientation is what makes the design extensible without violating the narrow-kernel principle, and is, more than any individual direction listed here, the contribution this work hopes to leave behind.
